\documentclass[11pt,a4paper,final]{article}
\usepackage[utf8]{inputenc}
\usepackage[portuguese]{babel}
\usepackage[T1]{fontenc}
\usepackage{amsmath}
\usepackage{amsfonts}
\usepackage{amssymb}
\usepackage{graphicx}
\usepackage{color}
\usepackage[left=2cm,right=2cm,top=1cm,bottom=2cm]{geometry}
\author{Tiago Á. Oliveira}
\title{Relatorio Individual}
\begin{document}
\thispagestyle{empty}
\begin{center}
\includegraphics[scale=1]{logo_UA}
\linebreak\linebreak
\LARGE\textbf{Eletricidade \& Magnetismo}
\normalsize\linebreak
(2013-2014)
\linebreak\linebreak
\textbf{Relat\'{o}rio Individual de Laborat\'{o}rio}
\linebreak
Turmas P3 e P10
\linebreak\linebreak
\underline{Informa\c{c}\~{o}es Relevantes}
\end{center}

\begin{itemize}
\item O relat\'{o}rio ser\'{a} realizado \underline{individualmente};

\item O local da realiza\c{c}\~{a}o do relat\'{o}rio ser\'{a} indicado oportunamente;

\item A data do relat\'{o}rio \'{e} a combinada anteriormente (ter\c{c}a-feira, 10 de dezembro de 2013), na hora habitual da aula (14h para a turma P3 e 16h para a turma P10);

\item Ter\'{a} a dura\c{c}\~{a}o exata de 2 horas, \underline{sem qualquer toler\^{a}ncia};

\item A mat\'{e}ria incidir\'{a} sobre a totalidade dos trabalhos pr\'{a}ticos, mas com forte incid\^{e}ncia nos dois \'{u}ltimos;

\item A distribui\c{c}\~{a}o, em termos de cota\c{c}\~{a}o final do relat\'{o}rio é 25\% para os primeiros quatro trabalhos e 75\% para os dois \'{u}ltimos;

\item O enunciado do relat\'{o}rio ter\'{a} 2 vers\~{o}es distintas mas de igual dificuldade. 

\underline{Dever\~{a}o indicar a vers\~{a}o (A ou B) no in\'{i}cio da folha de resposta};

\item Ser\~{a}o feitas perguntas e ordenadas realiza\c{c}\~{o}es que n\~{a}o correspondem necessariamente \`{a}s que constam nos guias;

\item \underline{Dever\~{a}o justificar todas as respostas};

\item Os dados experimentais a utilizar ser\~{a}o fornecidos;

\item Ser\'{a} pedida a realiza\c{c}\~{a}o de pelo menos 1 gr\'{a}fico em papel milim\'{e}trico ou Excel/MatLAB, ficando ao crit\'{e}rio de cada um;

\item N\~{a}o ser\'{a} permitida a consulta dos guias, fichas e dos registos efetuados aquando da realiza\c{c}\~{a}o dos trabalhos pr\'{a}ticos;

\item N\~{a}o haver\'{a} formul\'{a}rio para as leis fundamentais (Lei de Ohm, Kirschoff, Faraday, Lenz, etc.). Equa\c{c}\~{o}es de dedu\c{c}\~{a}o complexa/demorada, caso sejam necess\'{a}rias, estar\~{a}o presentes no enunciado;

\item Podem usar m\' {a}quina de calcular com capacidades gr\'{a}ficas;

\item O uso do computador/tablet est\'{a} restringido apenas para uso no tratamento de dados (Excel, MatLAB, etc.) n\~{a}o podendo ser usado para consulta de PDF's ou liga\c{c}\~{a}o \`{a} Internet.

\underline{A viola\c{c}\~{a}o desta norma implica anula\c{c}\~{a}o imediata da prova};

\item O telem\'{o}vel dever\'{a} estar desligado durante a prova;

\item Qualquer outra questão que não conste desta lista será decidida no momento.
\end{itemize}

\end{document}